
\section{磁盘结构与访问时间}
\concept{磁盘结构}
\begin{points}
  \item 磁盘由一个或多个盘片组成，盘片表面覆盖磁性材料。
  \item 每个盘面对应一个磁头，磁臂沿半径方向移动。
  \item 磁道是盘面上的同心圆，扇区是磁道上的固定大小区域。
  \item 柱面是所有盘面上半径相同的磁道集合。
  \item 一个物理块可用柱面号、磁头号、扇区号表示。
\end{points}
\plain{磁盘访问慢主要慢在机械移动：寻道和旋转。调度算法本质是少移动磁臂。}
\exam{常考按磁道请求序列计算移动磁道数，区分 SCAN/C-SCAN/LOOK/C-LOOK。}


\concept{磁盘访问时间}
\begin{points}
  \item 寻道时间：磁臂移动到目标柱面/磁道所需时间。
  \item 旋转延迟：目标扇区旋转到磁头下方所需时间，平均为旋转一周时间的一半。
  \item 传输时间：数据从磁盘读出或写入磁盘的时间。
  \item 总访问时间约为：寻道时间 + 旋转延迟 + 传输时间。
  \item 机械磁盘优化重点通常是减少寻道时间。
\end{points}
\plain{文件怎么读有两层：逻辑上记录如何组织，物理上块如何放在磁盘。}
\exam{常考顺序/直接访问，顺序文件、索引文件、索引顺序文件、散列文件适用场景。}


\section{磁盘调度算法}
\concept{FCFS}
\begin{points}
  \item 按磁盘请求到达顺序服务。
  \item 优点：公平、简单。
  \item 缺点：磁头来回移动可能很大，平均寻道距离长。
\end{points}
\plain{谁先来谁先上，最朴素但容易让短作业排在长作业后面。}
\exam{常考护航效应和平均等待时间计算。}


\concept{SSTF 最短寻道时间优先}
\begin{points}
  \item 每次选择距离当前磁头位置最近的请求。
  \item 优点：平均寻道距离通常较短。
  \item 缺点：远处请求可能长期等待，产生饥饿。
\end{points}
\plain{磁盘访问慢主要慢在机械移动：寻道和旋转。调度算法本质是少移动磁臂。}
\exam{常考按磁道请求序列计算移动磁道数，区分 SCAN/C-SCAN/LOOK/C-LOOK。}


\concept{SCAN 电梯算法}
\begin{points}
  \item 磁头像电梯一样沿一个方向移动，服务沿途请求，到达端点后反向。
  \item 优点：比 FCFS 稳定，避免频繁来回。
  \item 缺点：靠近两端的请求等待时间可能不同。
\end{points}
\plain{磁盘访问慢主要慢在机械移动：寻道和旋转。调度算法本质是少移动磁臂。}
\exam{常考按磁道请求序列计算移动磁道数，区分 SCAN/C-SCAN/LOOK/C-LOOK。}


\concept{C-SCAN 循环扫描}
\begin{points}
  \item 磁头只沿一个方向服务，到达一端后快速返回另一端，不在返回途中服务。
  \item 提供更均匀的等待时间。
  \item 适合请求分布较均匀的系统。
\end{points}
\plain{磁盘访问慢主要慢在机械移动：寻道和旋转。调度算法本质是少移动磁臂。}
\exam{常考按磁道请求序列计算移动磁道数，区分 SCAN/C-SCAN/LOOK/C-LOOK。}


\concept{LOOK 与 C-LOOK}
\begin{points}
  \item LOOK 不一定走到磁盘端点，只走到当前方向最远请求处就反向。
  \item C-LOOK 只在有请求的范围内循环扫描，从一端最远请求跳到另一端最远请求。
  \item 它们是 SCAN/C-SCAN 的实际优化版本，减少无意义移动。
\end{points}
\plain{磁盘访问慢主要慢在机械移动：寻道和旋转。调度算法本质是少移动磁臂。}
\exam{常考按磁道请求序列计算移动磁道数，区分 SCAN/C-SCAN/LOOK/C-LOOK。}


\section{磁盘管理与容错}
\concept{磁盘格式化}
\begin{points}
  \item 低级格式化把磁盘划分为扇区，为每个扇区建立控制信息。
  \item 分区把磁盘划分为一个或多个逻辑区域。
  \item 高级格式化/逻辑格式化在分区上建立文件系统数据结构，如空闲空间表、根目录、inode 区等。
\end{points}
\plain{磁盘访问慢主要慢在机械移动：寻道和旋转。调度算法本质是少移动磁臂。}
\exam{常考按磁道请求序列计算移动磁道数，区分 SCAN/C-SCAN/LOOK/C-LOOK。}


\concept{引导块与坏块}
\begin{points}
  \item 引导块保存启动操作系统所需的引导代码或引导信息。
  \item 坏块是无法可靠读写的磁盘块。
  \item 坏块管理可通过扇区备用、坏块链表、控制器重映射等方式处理。
\end{points}
\plain{磁盘管理不仅是读写，还包括初始化、启动、坏块处理和容错。}
\exam{常考低级/高级格式化、引导块作用、RAID 提高可靠性/性能的方式。}


\concept{交换空间管理}
\begin{points}
  \item 交换空间用于保存被换出的进程映像或页面。
  \item 交换空间可作为专门分区，也可作为文件系统中的交换文件。
  \item 专门分区管理简单且效率高；交换文件更灵活但可能受文件系统开销影响。
\end{points}
\plain{把这个知识点放回本章主线理解：它回答的是 OS 为了管理资源、隐藏硬件、保证并发正确性而采用的某个对象、规则或步骤。}
\exam{常考选择/填空/简答，让你判断概念归属、比较相近概念，或按“定义-作用-机制-易错点”展开。}


\concept{RAID 磁盘容错}
\begin{points}
  \item RAID 通过磁盘冗余提高性能和可靠性。
  \item RAID 0 条带化提高性能但无冗余。
  \item RAID 1 镜像保存完整副本，可靠性高但空间开销大。
  \item RAID 5 使用分布式奇偶校验，兼顾容量和容错。
  \item RAID 6 可容忍两块磁盘故障，可靠性更高。
\end{points}
\plain{磁盘访问慢主要慢在机械移动：寻道和旋转。调度算法本质是少移动磁臂。}
\exam{常考按磁道请求序列计算移动磁道数，区分 SCAN/C-SCAN/LOOK/C-LOOK。}


